\documentclass[11pt,a4paper]{article}
\usepackage[T1]{fontenc}
\usepackage{lmodern,amsmath,amssymb,amsthm,mathtools,booktabs,microtype}
\usepackage[margin=27mm]{geometry}
\usepackage[colorlinks=true,linkcolor=blue,citecolor=blue,urlcolor=blue]{hyperref}
\setlength{\emergencystretch}{2em}
\newtheorem{theorem}{Theorem}[section]
\newtheorem{lemma}[theorem]{Lemma}
\newtheorem{proposition}[theorem]{Proposition}
\newtheorem{corollary}[theorem]{Corollary}
\theoremstyle{definition}\newtheorem{definition}[theorem]{Definition}
\newcommand{\Q}{\mathbb Q}\newcommand{\Z}{\mathbb Z}
\title{Obstructions to automatic first conics in pole certificates}
\author{}\date{}
\begin{document}\maketitle
\begin{abstract}
Writing a pole input as $b=(1-e^2)r^2$ makes the first congruent-number
conic automatic. This does not ensure that the second pole equations have
rational witnesses. We give two dense families of prepared counterexamples,
including arbitrarily small positive inputs. One has the fixed value
$e=9/41$ and square inputs. The other occurs on a rational contact map for
which $e,1-e,1+e$ all have completely split numerators. The obstructions use
exact nonzero coefficients in Tunnell's theorem and allow a zero second
ordinate. They concern prepared multipliers; no assertion is made that the
specified multipliers occur in a complete integer-defining source.
\end{abstract}

\section{The pole predicate and noncongruent areas}

The congruent-number pole equations of \cite{Pole} are
\begin{equation}\label{eq:pole}
 e^2+bp^2=1,\qquad uv=e(1-b),\qquad u^2-v^2=bq^2.
\end{equation}
A \emph{prepared} input in this note means a rational number $b$ with
square reduced denominator satisfying
\begin{equation}\label{eq:prepared}
                 0<b<1,\qquad b\in\Q_2^{\times2},\qquad v_2(b)\ge3.
\end{equation}
If $b=(1-e^2)r^2$ with $r\ne0$, then $p=1/r$ satisfies the first
equation. The remaining equations need a separate rational realization.

\begin{lemma}\label{lem:area}
Suppose $0<b<1$, $e>0$, and put $A=e(1-b)$. If the last two equations
of \eqref{eq:pole} have a rational solution with $q\ne0$, then $bA$ is
the area of a rational right triangle. If $A$ is not a rational square,
there is no solution with $q=0$.
\end{lemma}
\begin{proof}
For $q\ne0$, the positive rational numbers
\[
 \left|\frac{2A}{q}\right|,\qquad |bq|,\qquad
 \left|\frac{u^2+v^2}{q}\right|
\]
are the sides of a right triangle. Indeed
$(2uv)^2+(u^2-v^2)^2=(u^2+v^2)^2$, and its area is $bA$.
If $q=0$, then $u^2=v^2$. Since $uv=A>0$, one has $u=v$ and
$A=u^2$, proving the second assertion.
\end{proof}

The property of being the area of a rational right triangle is unchanged
by multiplication by a nonzero rational square. We will use the following
two noncongruent integers.

\begin{proposition}\label{prop:noncongruent}
Neither $17$ nor
\[
                    N_*=15\,499\,155=3\cdot5\cdot7\cdot17\cdot19\cdot457
\]
is the area of a rational right triangle.
\end{proposition}
\begin{proof}
Tunnell's introduction theorem, part~(a), states that $a(n)\ne0$ excludes
$n$ as such an area, where
\begin{equation}\label{eq:tunnell}
 \sum_{n\ge1}a(n)z^n
 =z\prod_{m\ge1}(1-z^{8m})(1-z^{16m})
       \sum_{c\in\Z}z^{2c^2}.
\end{equation}
This is the unconditional necessary direction of the theorem
\cite[p.~323]{Tunnell}; its converse is not used.

The first product, including its initial $z$, is
$z-z^9-2z^{17}+O(z^{25})$, and the theta series is
$1+2z^2+2z^8+O(z^{18})$. Thus $a(17)=-4$.

For the larger integer, Euler's pentagonal identity gives the finite sum
\begin{equation}\label{eq:finite-sum}
 a(n)=\sum_{\substack{j,k,c\in\Z\\
      4j(3j-1)+8k(3k-1)+2c^2=n-1}}(-1)^{j+k}.
\end{equation}
The exact verifier enumerates the $2273$ possible first exponents and
$1607$ possible second exponents at $n=N_*$. It determines $c$ by exact
integer square-root testing. There are $952$ exponent pairs contributing
$+2$ and $1080$ contributing $-2$, with no $c=0$ contribution. Hence
\[
                        a(N_*)=2(952-1080)=-256.
\]
As an independent finite check, the numbers of integer solutions of
$2x^2+y^2+8z^2=N_*$ and $2x^2+y^2+32z^2=N_*$ are $10752$ and
$5120$, respectively; their difference $5120-10752/2$ is again $-256$.
The corresponding counts for $17$ are $16$ and $4$. The nonzero
coefficients in \eqref{eq:tunnell} prove both assertions.
\end{proof}

\section{A fixed conic with square prepared inputs}

\begin{theorem}\label{thm:small}
Fix
\[
                 e=\frac9{41},\qquad R=1-e^2=\left(\frac{40}{41}\right)^2.
\]
There is a dense subset of $(0,1)$ consisting of prepared rational squares
$b=Rr^2$ for which the first equation of \eqref{eq:pole} has the explicit
witness $p=1/r$, but the last two equations have no rational solution,
even with $q=0$ allowed. One example is
\[
 b=\left(\frac{348}{349}\right)^2,\qquad
 r=\frac{3567}{3490},\qquad p=\frac{3490}{3567}.
\]
\end{theorem}
\begin{proof}
Consider the rational conic
\begin{equation}\label{eq:smallconic}
 t^2+697w^2=1,\qquad
 t=\frac{1-697z^2}{1+697z^2},\qquad
 w=\frac{2z}{1+697z^2}.
\end{equation}
For nonzero rational $z$, neither $t$ nor $w$ vanishes: $697=17\cdot41$
is not a rational square. Put
\[
                b=t^2,\qquad r=\frac{41t}{40},\qquad p=\frac{40}{41t}.
\]
Then $0<b<1$, $b=Rr^2$, and $e^2+bp^2=1$. If $z$ is a $2$-adic
unit, then $697z^2\equiv1\pmod8$. The numerator of $t$ has valuation
at least three and its denominator has valuation one. Hence
$v_2(t)\ge2$, so $b$ satisfies \eqref{eq:prepared}. Its reduced
denominator is a square because $b=t^2$.

The second-pole factor is
\[
                     A=e(1-b)=\frac9{41}\,697w^2=17(3w)^2.
\]
It is not a rational square, so $q=0$ is excluded by Lemma~\ref{lem:area}.
If $q\ne0$, the same lemma would make $bA=17(3tw)^2$ a rational
right-triangle area. This contradicts Proposition~\ref{prop:noncongruent}.

Rational $2$-adic units are dense in $\mathbb R$: fractions with odd
numerator and arbitrarily large odd denominator approximate every real
number. On $0<z<1/\sqrt{697}$, the continuous function $t(z)^2$
has image $(0,1)$. This proves density and gives arbitrarily small prepared
counterexamples. At $z=1/697$, one has $(t,w)=(348/349,1/349)$,
which gives the displayed example and $v_2(b)=4$.
\end{proof}

The fixed first conic is not itself inconsistent with the full pole
predicate. At the same $e=9/41$, the prepared input $b=R$ has the full
witness
\[
                         p=1,\qquad u=e,\qquad v=e^2,\qquad q=e.
\]
The obstruction depends on the prepared multiplier.

\section{A completely split six-contact map}

Consider the rational functions
\begin{equation}\label{eq:split}
 e(\delta)=\frac{49\delta^2-1}{120\delta^3},\qquad
 R(\delta)=\frac{\displaystyle\prod_{a\in\{3,5,8\}}
                    (\delta^2-a^{-2})}{\delta^6}.
\end{equation}
The identities
\begin{equation}\label{eq:splitids}
 \begin{split}
 e(\delta)&=\frac{(7\delta-1)(7\delta+1)}{120\delta^3},\\
 1-e(\delta)&=\frac{(3\delta-1)(5\delta-1)(8\delta+1)}{120\delta^3},\\
 1+e(\delta)&=\frac{(3\delta+1)(5\delta+1)(8\delta-1)}{120\delta^3}
 \end{split}
\end{equation}
show that $R=1-e^2$ and that all the displayed numerators split over $\Q$.
Thus $b=R(\delta)r^2$ and $p=1/r$ make the first pole conic automatic
using only rational contact functions.

\begin{theorem}\label{thm:split}
At $\delta=7$, there is a dense subset of $(0,1)$ consisting of prepared
inputs $b=R(7)r^2$ for which the last two equations of \eqref{eq:pole}
have no rational solution, including solutions with $q=0$. In particular,
put
\[
 t_* =\frac{38864}{1235629},\qquad
 w_* =\frac{11759}{1235629},\qquad b_*=969t_*^2.
\]
Then $v_2(b_*)=8$, and the first conic has
\[
 e=\frac{20}{343},\qquad R(7)=969\left(\frac{11}{343}\right)^2,
 \qquad r=\frac{343t_*}{11},\qquad p=\frac{11}{343t_*}.
\]
\end{theorem}
\begin{proof}
Set $\beta=969=3\cdot17\cdot19$ and $\theta=457$. The conic
\begin{equation}\label{eq:largeconic}
                         \beta t^2+\theta w^2=1
\end{equation}
has the rational parametrization
\begin{equation}\label{eq:largeparam}
 \begin{split}
 t(z)&=\frac{8721+73120z-4113z^2}{1733(969+457z^2)},\\
 w(z)&=\frac{77520-17442z-36560z^2}{1733(969+457z^2)}.
 \end{split}
\end{equation}
It comes from lines through $(9/1733,80/1733)$, followed by reflection
in the $w$-axis. Direct substitution proves \eqref{eq:largeconic}.
At $z=1$ it gives $(t_*,w_*)$.

For every rational $z$, both coordinates are nonzero: a zero coordinate
would make $\beta$ or $\theta$ a rational square. Therefore
$b=\beta t(z)^2$ lies in $(0,1)$. Its reduced denominator is a square.
Indeed a negative valuation of $t$ at a prime dividing $\beta$ would
give $\beta t^2$ negative odd valuation, whereas $\theta w^2$ has even
valuation because $\gcd(\beta,\theta)=1$. These two valuations cannot
cancel to give~$1$. Thus $t$ is integral at all primes dividing $\beta$;
at every other prime a negative valuation of $b$ is even.

Suppose now that $z$ is a $2$-adic unit. Modulo $8$, the numerator of
$t(z)$ is $1-z^2=0$, while its denominator has valuation one. Thus
$v_2(t)\ge2$. Since $\beta\equiv1\pmod8$, the value $b=\beta t^2$
is a $2$-adic square and is prepared. The value at $z=1$ has
$v_2(t_*)=4$ and $v_2(b_*)=8$.

The first-conic identity follows from \eqref{eq:splitids}. For the
remaining equations note that $e=20/343=35(2/49)^2$ and
$1-b=457w^2$. Consequently
\begin{equation}\label{eq:areas}
 A=e(1-b)=15995\left(\frac{2w}{49}\right)^2,
 \qquad
 bA=N_*\left(\frac{2tw}{49}\right)^2.
\end{equation}
The first factor $15995=5\cdot7\cdot457$ is not a square, so $q=0$
is impossible. A nonzero $q$ would give a rational right triangle of area
$bA$ by Lemma~\ref{lem:area}, contradicting the noncongruence of $N_*$.

Finally the rational parametrization \eqref{eq:largeparam} covers the
real conic except at most one point. The function $\beta t^2$ has
image $[0,1]$ on that conic. Every value in $(0,1)$ has a finite real
parameter preimage, after choosing a suitable sign for $t$ or $w$.
Approximate that parameter by rational $2$-adic units. The resulting
prepared counterexamples are dense in $(0,1)$.
\end{proof}

Theorems~\ref{thm:small} and~\ref{thm:split} exclude a realization
theorem for \emph{every} prepared multiplier after the first conic has
been forced in the indicated ways. They do not exclude a more restricted
arithmetic source whose detector values avoid these families. Such a source
must prove the remaining rational realization for its actual multipliers;
the automatic first-conic identity and real smallness of the pole input
do not provide it.

\section{Exact verification}

The verifier
\nolinkurl{computation/pole_parameters/automatic_conic_v41/check.py}
checks the contact factorizations, both conic parametrizations, all displayed
specializations, preparation of the explicit examples, and the square
factors in \eqref{eq:areas}. It exhausts the four unit residue classes
modulo $8$ in the second family's valuation argument. The direct formal
product is checked through degree $257$ against the pentagonal enumeration.
For $17$ and $N_*$, independent exhaustive ternary counts agree with the
computed coefficients. The saved certificate includes the signed-sum
counts, exact example coordinates, a positive-control witness, and the
source hash of the verifier. No rank computation or conjectural converse
to Tunnell's theorem is used.

\begin{thebibliography}{9}
\bibitem{Pole}
\emph{Pole equations from simultaneous congruent-number twists},
\nolinkurl{papers/pole_parameters/congruent_twist_poles_20261114_v27.tex}.
\bibitem{Tunnell}
J. B. Tunnell, \emph{A Classical Diophantine Problem and Modular Forms
of Weight $3/2$}, Invent. Math. \textbf{72} (1983), 323--334.
\href{https://doi.org/10.1007/BF01389327}{doi:10.1007/BF01389327}.
Primary scan:
\url{https://sites.math.rutgers.edu/~zeilberg/EM22/JT1983.pdf}.
\end{thebibliography}
\end{document}
